\documentclass[11pt]{article}

% Page geometry
\usepackage[margin=1in]{geometry}

% Math
\usepackage{amsmath,amssymb,amsthm}

% Tables
\usepackage{booktabs}
\usepackage{multirow}
\usepackage{array}

% Figures and graphics
\usepackage{graphicx}
\usepackage{xcolor}

% References and links
\usepackage[colorlinks=true,linkcolor=blue,citecolor=blue,urlcolor=blue]{hyperref}
\usepackage{natbib}
\bibliographystyle{plainnat}

% Formatting
\usepackage{microtype}
\usepackage{enumitem}

% Theorems
\newtheorem{theorem}{Theorem}
\newtheorem{proposition}{Proposition}
\newtheorem{lemma}{Lemma}

\title{Lossless 12$\times$ KV Cache Compression in Hybrid Attention Models\\via Combined Entropy-Adaptive Eviction and Vector Quantization}

\author{
SCJedi\\
Independent Researcher\\
\url{https://github.com/SCJedi}
}

\date{March 2026}

\begin{document}

\maketitle

%==============================================================================
\begin{abstract}
%==============================================================================
The key-value (KV) cache is the dominant memory bottleneck for long-context large language model inference.  We present a combined compression pipeline that applies entropy-adaptive token eviction followed by TurboQuant-style vector quantization, exploiting the orthogonality of these two techniques: eviction reduces the number of cached tokens while quantization reduces the bits per token, yielding multiplicative memory savings.  We evaluate this pipeline across four models---GPT-2 (124M), Qwen2.5-7B, Qwen3.5-4B, and Qwen3.5-9B---spanning standard dense attention, standard attention with grouped query attention (GQA), and hybrid attention-linear architectures.  On the Qwen3.5 family, which interleaves 8 full-attention layers with 24 Gated Delta Net linear-attention layers, the combined pipeline achieves \textbf{12$\times$ KV cache compression with zero generation quality loss}: BLEU 1.000, 100\% token match, and perplexity within 1.2--3.3\% of the uncompressed baseline.  On GPT-2, we demonstrate super-additive synergy, where combining 2$\times$ eviction with 4-bit quantization at 8$\times$ total compression outperforms standalone 4-bit quantization at only 4$\times$ compression.  We show that hybrid attention architectures are the critical enabler for lossless high-ratio compression: the linear-attention layers act as an error-correcting backbone that absorbs quantization noise before it propagates.  At 12$\times$ compression, Qwen3.5-9B's KV cache for a 128K-token context shrinks from 2.0\,GB to 171\,MB, enabling long-context inference on consumer GPUs.
\end{abstract}

%==============================================================================
\section{Introduction}
\label{sec:intro}
%==============================================================================

Large language models (LLMs) store key-value (KV) pairs at every layer during autoregressive generation, and this cache grows linearly with both sequence length and model depth.  For a 70B-parameter model at 4K context, the KV cache alone consumes approximately 10\,GB---often exceeding the model weights themselves in memory-constrained deployments.  This bottleneck limits batch size, context length, and deployment on consumer hardware.

Two families of techniques have emerged to compress the KV cache.  \emph{Token eviction} methods~\citep{zhang2024h2o,liu2024scissorhands,xiao2024efficient} reduce memory by discarding tokens deemed unimportant, keeping only a subset of the cached entries.  \emph{Vector quantization} methods~\citep{hooper2024kvquant,kang2024gear,liu2024kivi,duanmu2024turboquant} reduce memory by representing each cached vector at lower bit-width, typically 2--4 bits instead of 16.  These two approaches operate on orthogonal axes---eviction reduces the \emph{number} of vectors while quantization reduces the \emph{bits per vector}---suggesting that they should compose multiplicatively.

In this work, we demonstrate that:
\begin{enumerate}[leftmargin=*]
    \item Combined eviction and quantization achieves exactly multiplicative memory reduction and exhibits \textbf{super-additive quality retention}: the combined quality loss is less than the sum of individual losses, because eviction removes the tokens that would contribute the most quantization noise.
    \item \textbf{Model architecture determines compression robustness.}  Hybrid attention models (mixing full attention with linear attention) achieve lossless compression at ratios where standard transformers degrade substantially.
    \item Qwen3.5-4B and Qwen3.5-9B sustain \textbf{BLEU 1.000 and 100\% token match at 12$\times$ compression}, with perplexity ratios of 1.033 and 1.012 respectively---effectively lossless.
    \item The enabling mechanism is a combination of four architectural properties: an error-correcting linear-attention backbone, grouped query attention (GQA), large head dimensions, and narrow attention entropy distributions.
\end{enumerate}

%==============================================================================
\section{Related Work}
\label{sec:related}
%==============================================================================

\paragraph{KV Cache Eviction.}
H2O~\citep{zhang2024h2o} identifies ``heavy hitter'' tokens via cumulative attention scores and retains them with a recent-token window.  Scissorhands~\citep{liu2024scissorhands} exploits attention persistence---tokens important at one step tend to remain important.  StreamingLLM~\citep{xiao2024efficient} demonstrates that retaining a small set of initial ``attention sink'' tokens plus a sliding window suffices for perplexity maintenance.  Our entropy-adaptive approach generalizes these by assigning per-head retention budgets proportional to measured attention entropy, which naturally handles the heterogeneous head behaviors observed in practice.

\paragraph{KV Cache Quantization.}
KIVI~\citep{liu2024kivi} applies per-channel quantization to keys and per-token quantization to values, achieving 2-bit KV caches with minimal quality loss on specific models.  KVQuant~\citep{hooper2024kvquant} adds per-channel key quantization, non-uniform quantization, and outlier handling.  Gear~\citep{kang2024gear} uses low-rank approximation plus quantization of residuals.  TurboQuant~\citep{duanmu2024turboquant} applies random Hadamard rotation before Lloyd-Max scalar quantization with QJL residual correction, which we adopt as our quantization backbone.

\paragraph{Hybrid Attention Architectures.}
Jamba~\citep{lieber2024jamba} interleaves Mamba~\citep{gu2024mamba} state-space layers with transformer attention layers at a 7:1 ratio.  Qwen3.5~\citep{qwen2025qwen3} uses a similar hybrid design, replacing Mamba with Gated Delta Net linear attention: 8 full-attention layers interleaved with 24 linear-attention layers in a 32-layer model.  These architectures inherently reduce KV cache size since only the full-attention layers maintain a traditional growing cache.

\paragraph{Combined Approaches.}
To our knowledge, this is the first systematic study of combining token eviction with vector quantization for KV cache compression, and the first to demonstrate that hybrid attention architectures are the critical enabler for lossless high-ratio combined compression.

%==============================================================================
\section{Method}
\label{sec:method}
%==============================================================================

%----------------------------------------------------------------------
\subsection{Entropy-Adaptive Token Eviction}
\label{sec:eviction}
%----------------------------------------------------------------------

We assign each attention head a KV cache retention budget proportional to its measured attention entropy.  Given a calibration set of sequences, we compute the Shannon entropy of each head's attention distribution:
\begin{equation}
    H(h) = -\sum_{i=1}^{n} a_i^{(h)} \log_2 a_i^{(h)}
    \label{eq:entropy}
\end{equation}
where $a_i^{(h)}$ are the softmax attention weights for head $h$.  Low-entropy (``sink'') heads concentrate attention on one or two positions and require very few KV entries.  High-entropy (``diffuse'') heads attend broadly and require more entries.  The retention budget for head $h$ is:
\begin{equation}
    r_h = r_{\min} + (r_{\max} - r_{\min}) \cdot \frac{H(h) - H_{\min}}{H_{\max} - H_{\min}}
    \label{eq:budget}
\end{equation}
where $r_{\min}$ and $r_{\max}$ bound the per-head keep ratio for a target overall compression $R_e$.

\paragraph{GQA-Aware Min-Entropy Bottleneck.}
In models with grouped query attention (GQA), multiple query heads share a single KV head.  Evicting a token from the shared KV cache affects all query heads in the group.  We set the retention budget for KV group $g$ using the minimum entropy across its query heads:
\begin{equation}
    r_g = f\!\left(\min_{j \in \mathrm{group}(g)} H_j\right)
    \label{eq:gqa_bottleneck}
\end{equation}
This ensures the most focused (lowest-entropy) query head in each group is never under-provisioned.

%----------------------------------------------------------------------
\subsection{TurboQuant Vector Quantization}
\label{sec:turboquant}
%----------------------------------------------------------------------

We adopt the TurboQuant pipeline~\citep{duanmu2024turboquant} for quantizing the surviving KV vectors:

\begin{enumerate}[leftmargin=*]
    \item \textbf{Random Hadamard rotation.}  A fixed random orthogonal matrix $R$ is applied to each value vector: $\tilde{v} = Rv$.  This decorrelates dimensions and makes the marginal distribution per coordinate approximately Gaussian (by the CLT), which is the optimal input distribution for scalar quantization.
    \item \textbf{Lloyd-Max scalar quantization.}  Each rotated coordinate is independently quantized to $b$ bits using a Lloyd-Max codebook~\citep{lloyd1982least,max1960quantizing}, which minimizes mean squared error for a given bit budget and input distribution.
    \item \textbf{QJL residual correction.}  A Quantized Johnson-Lindenstrauss (QJL) sketch captures residual quantization error, providing a probabilistic correction at inference time.
\end{enumerate}

The quantization compression ratio is $R_q = b_0 / b$, where $b_0 = 16$ is the original bit-width.  At 4-bit quantization, $R_q = 4$.

\paragraph{Per-Vector L2 Normalization.}
Production models with RoPE~\citep{su2024roformer} positional encodings exhibit significant magnitude variation across key vectors at different sequence positions.  Without normalization, large-magnitude outliers dominate the Lloyd-Max codebook, wasting bits on range coverage.  We apply per-vector L2 normalization before quantization and store the scalar norm separately, denormalizing after dequantization.  This was critical for Qwen2.5-7B and the Qwen3.5 family.

%----------------------------------------------------------------------
\subsection{Combined Pipeline: Evict First, Then Quantize}
\label{sec:combined}
%----------------------------------------------------------------------

The combined pipeline applies eviction before quantization:
\begin{enumerate}[leftmargin=*]
    \item Compute full attention weights using the uncompressed KV cache.
    \item Apply entropy-adaptive eviction, retaining $n/R_e$ tokens per head.
    \item Apply TurboQuant quantization to the surviving vectors at $b$ bits.
\end{enumerate}

The combined memory per head is:
\begin{equation}
    M_{\text{combined}} = \frac{n}{R_e} \cdot d \cdot \frac{b_0}{R_q} = \frac{M_0}{R_e \cdot R_q}
    \label{eq:combined_memory}
\end{equation}
where $d$ is the head dimension.  \textbf{The combined compression ratio is exactly $R_e \times R_q$.}  This is an arithmetic identity---the two techniques operate on orthogonal axes (token count vs.\ bits per vector) with no memory overhead from combining them.

The order of operations matters.  Eviction before quantization is strictly preferable because: (1) eviction decisions depend on attention weights computed from full-precision keys, so quantizing first would degrade eviction quality; (2) quantizing fewer surviving vectors is computationally cheaper; and (3) Lloyd-Max calibration on the post-eviction distribution better matches the actual data being quantized.

%----------------------------------------------------------------------
\subsection{Theoretical Analysis}
\label{sec:theory}
%----------------------------------------------------------------------

\paragraph{Quality Interaction Term.}
Let $A$ be the full attention matrix (post-softmax) and $V$ the full value matrix, so the output is $O = AV$.  Eviction replaces $A$ with $\tilde{A}$ (masked and renormalized); quantization replaces $V$ with $\hat{V}$.  The combined output error decomposes as:
\begin{equation}
    \tilde{A}\hat{V} - AV = \underbrace{(\tilde{A} - A)V}_{\text{eviction error}} + \underbrace{A(\hat{V} - V)}_{\text{quantization error}} + \underbrace{(\tilde{A} - A)(\hat{V} - V)}_{\text{interaction}}
    \label{eq:error_decomp}
\end{equation}

The interaction term is bounded by:
\begin{equation}
    \|(\tilde{A} - A)(\hat{V} - V)\|_F \leq \|\tilde{A} - A\|_F \cdot \|\hat{V} - V\|_F = O(\epsilon_e \cdot \epsilon_q)
    \label{eq:interaction_bound}
\end{equation}
where $\epsilon_e$ and $\epsilon_q$ are the individual error magnitudes.  When both perturbations are small (e.g., 5\% each), the interaction is $\sim$0.25\%---negligible.  Empirically, we observe this interaction is often \emph{negative} (synergistic benefit), because eviction removes the tokens that contribute the most quantization noise.

\paragraph{Quantization Error Scaling with Head Entropy.}
For head $h$ with attention weights $a_i^{(h)}$ and i.i.d.\ quantization errors $\delta_i$ with $\mathbb{E}[\|\delta_i\|^2] = \sigma_q^2 d$, the expected output error is:
\begin{equation}
    \mathbb{E}[\|\epsilon_h\|^2] = \sigma_q^2 \cdot d \cdot \sum_{i=1}^{n} \left(a_i^{(h)}\right)^2 = \sigma_q^2 \cdot d \cdot 2^{-H_2(h)}
    \label{eq:quant_error_entropy}
\end{equation}
where $H_2(h) = -\log_2 \sum_i (a_i^{(h)})^2$ is the R\'{e}nyi entropy of order 2.  \textbf{High-entropy heads are more tolerant of quantization noise}: their diffuse attention averages over many value vectors, suppressing per-vector quantization error by a law-of-large-numbers effect.

\paragraph{Optimal Per-Head Bit Allocation.}
Minimizing total output error across $H$ heads subject to a total bit budget $B_{\text{total}} = \sum_h b_h$, and using the standard rate-distortion bound $\sigma_q^2(b) \approx c \cdot 2^{-2b}$, Lagrange optimization yields:
\begin{equation}
    \boxed{b_h^* = \bar{b} + \frac{\bar{H}_2 - H_2(h)}{2}}
    \label{eq:optimal_bits}
\end{equation}
where $\bar{b} = B_{\text{total}}/H$ and $\bar{H}_2$ is the mean R\'{e}nyi entropy.  This allocates \emph{more} bits to low-entropy (focused) heads and \emph{fewer} bits to high-entropy (diffuse) heads---each bit of R\'{e}nyi entropy difference corresponds to a half-bit reallocation.

%==============================================================================
\section{Experimental Setup}
\label{sec:setup}
%==============================================================================

%----------------------------------------------------------------------
\subsection{Models}
%----------------------------------------------------------------------

We evaluate four models spanning three architectural classes (Table~\ref{tab:models}).

\begin{table}[htbp]
\centering
\caption{Model architectures.  ``Attn layers'' counts full softmax-attention layers (with traditional KV cache); remaining layers use linear attention (Qwen3.5) or are all attention (GPT-2, Qwen2.5).  GQA ratio is the number of query heads per KV head.}
\label{tab:models}
\begin{tabular}{@{}lcccccc@{}}
\toprule
Model & Params & Layers & Attn layers & GQA ratio & head\_dim & Architecture \\
\midrule
GPT-2 & 124M & 12 & 12 & 1:1 & 64 & Standard MHA \\
Qwen2.5-7B & 7B & 28 & 28 & 7:1 & 128 & Standard + GQA \\
Qwen3.5-4B & 4B & 32 & 8 & 4:1 & 256 & Hybrid + GQA \\
Qwen3.5-9B & 9B & 32 & 8 & 4:1 & 256 & Hybrid + GQA \\
\bottomrule
\end{tabular}
\end{table}

GPT-2~\citep{radford2019language} serves as a baseline with standard multi-head attention (MHA) and no GQA.  Qwen2.5-7B~\citep{qwen2024qwen25} uses standard attention across all 28 layers with 7:1 GQA.  The Qwen3.5 models~\citep{qwen2025qwen3} use a hybrid architecture: 8 full-attention layers (at positions 3, 7, 11, 15, 19, 23, 27, 31) interleaved with 24 Gated Delta Net linear-attention layers, with 4:1 GQA and head\_dim${}=256$.

%----------------------------------------------------------------------
\subsection{Configurations}
%----------------------------------------------------------------------

Each model is evaluated on up to 10 compression configurations:
\begin{itemize}[leftmargin=*]
    \item \textbf{Baseline:} Full KV cache at 16-bit precision.
    \item \textbf{Eviction only:} Entropy-adaptive eviction at 2$\times$ and 3$\times$ compression.
    \item \textbf{Quantization only:} TurboQuant at 5-bit (3.2$\times$), 4-bit (4$\times$), and 3-bit (5.3$\times$).
    \item \textbf{Combined:} Eviction 2$\times$ + quant 5-bit (6.4$\times$), eviction 2$\times$ + quant 4-bit (8$\times$), eviction 3$\times$ + quant 4-bit (12$\times$).
    \item \textbf{Entropy-adaptive quantization:} Per-head bit allocation with low-entropy heads receiving more bits.
\end{itemize}

%----------------------------------------------------------------------
\subsection{Metrics}
%----------------------------------------------------------------------

All metrics are computed against the full-cache baseline generation:
\begin{itemize}[leftmargin=*]
    \item \textbf{BLEU}~\citep{papineni2002bleu}: $n$-gram overlap with baseline generation (normalized so baseline = 1.0).
    \item \textbf{ROUGE-L}~\citep{lin2004rouge}: Longest common subsequence overlap.
    \item \textbf{Token match rate}: Fraction of generated tokens identical to baseline.
    \item \textbf{Perplexity ratio}: $\text{PPL}_{\text{compressed}} / \text{PPL}_{\text{baseline}}$ (1.0 = no degradation).
\end{itemize}

%----------------------------------------------------------------------
\subsection{Evaluation Protocol}
%----------------------------------------------------------------------

\begin{itemize}[leftmargin=*]
    \item \textbf{GPT-2:} 8 WikiText-2 validation sequences, 128-token context, 20-token greedy generation.  CPU inference (FP32 weights).
    \item \textbf{Qwen models:} 6 sequences, 256-token context, 30-token greedy generation.  GPU inference (RTX 3060 12\,GB; Qwen2.5-7B and Qwen3.5 at Q4\_K\_M GGUF weight quantization).
    \item All experiments use seed 42 for reproducibility.
\end{itemize}

%==============================================================================
\section{Results}
\label{sec:results}
%==============================================================================

%----------------------------------------------------------------------
\subsection{GPT-2: Super-Additive Synergy}
\label{sec:gpt2_results}
%----------------------------------------------------------------------

Table~\ref{tab:gpt2_results} shows results for GPT-2 across 12 configurations.

\begin{table}[htbp]
\centering
\caption{GPT-2 (124M) compression results.  The combined 2$\times$+4-bit configuration at 8$\times$ compression outperforms standalone 4-bit quantization at only 4$\times$ compression, demonstrating super-additive synergy.}
\label{tab:gpt2_results}
\begin{tabular}{@{}llccccc@{}}
\toprule
Configuration & Mode & $R$ & BLEU & ROUGE-L & Match & PPL ratio \\
\midrule
Full cache & Baseline & 1.0$\times$ & 1.000 & 1.000 & 100.0\% & 1.000 \\
\midrule
Eviction 2$\times$ & Eviction & 2.0$\times$ & 0.811 & 0.856 & 81.3\% & 1.027 \\
Eviction 4$\times$ & Eviction & 4.0$\times$ & 0.690 & 0.775 & 61.9\% & 1.121 \\
\midrule
Quant 4-bit & Quantization & 4.0$\times$ & 0.516 & 0.681 & 46.3\% & 1.089 \\
Quant 3-bit & Quantization & 5.3$\times$ & 0.358 & 0.525 & 24.4\% & 1.274 \\
Quant 2-bit & Quantization & 8.0$\times$ & 0.115 & 0.181 & 4.4\% & 3.643 \\
\midrule
\textbf{Combined 2$\times$+4bit} & \textbf{Combined} & \textbf{8.0$\times$} & \textbf{0.596} & \textbf{0.694} & \textbf{52.5\%} & \textbf{1.082} \\
Combined 2$\times$+3bit & Combined & 10.7$\times$ & 0.247 & 0.419 & 18.1\% & 1.344 \\
Combined 3$\times$+4bit & Combined & 12.0$\times$ & 0.464 & 0.581 & 34.4\% & 1.159 \\
Combined 3$\times$+3bit & Combined & 16.0$\times$ & 0.265 & 0.400 & 13.1\% & 1.465 \\
\midrule
Entropy quant (low-H $\to$ fewer bits) & Entropy & 5.3$\times$ & 0.316 & 0.475 & 18.8\% & 1.508 \\
Entropy quant (low-H $\to$ more bits) & Entropy & 5.3$\times$ & 0.205 & 0.369 & 13.1\% & 1.602 \\
\bottomrule
\end{tabular}
\end{table}

The key finding is \textbf{super-additive synergy} at 8$\times$ compression.  Comparing methods at equal compression:
\begin{itemize}[leftmargin=*]
    \item \textbf{Pure 2-bit quantization} at 8$\times$: BLEU 0.115, 4.4\% match, PPL ratio 3.64.
    \item \textbf{Combined 2$\times$+4-bit} at 8$\times$: BLEU 0.596, 52.5\% match, PPL ratio 1.08.
\end{itemize}
The combined approach achieves 5.2$\times$ higher BLEU at the same compression.  More remarkably, the combined 8$\times$ configuration (BLEU 0.596) \emph{outperforms} standalone 4-bit quantization at only 4$\times$ compression (BLEU 0.516).  Adding 2$\times$ eviction on top of 4-bit quantization simultaneously doubles compression and \emph{improves} quality---eviction removes the tokens that contribute the most quantization noise.

\paragraph{Entropy-Informed Bit Allocation.}
Allocating fewer bits to high-entropy heads (BLEU 0.316) outperforms the reverse allocation (BLEU 0.205) by 54\%, confirming the theoretical prediction from Equation~\ref{eq:quant_error_entropy} that diffuse heads tolerate more quantization noise.  However, neither entropy-informed strategy outperforms uniform 3-bit quantization (BLEU 0.358), suggesting the per-head reallocation magnitude exceeded the theoretical optimum of half a bit per unit of R\'{e}nyi entropy.

%----------------------------------------------------------------------
\subsection{Qwen3.5 Family: Lossless 12$\times$ Compression}
\label{sec:qwen35_results}
%----------------------------------------------------------------------

Tables~\ref{tab:qwen35_4b} and~\ref{tab:qwen35_9b} present results for the Qwen3.5 hybrid models.

\begin{table}[htbp]
\centering
\caption{Qwen3.5-4B compression results.  All 10 configurations achieve BLEU 1.000 and 100\% token match.  Maximum perplexity degradation at 12$\times$: 3.3\%.}
\label{tab:qwen35_4b}
\begin{tabular}{@{}llcccc@{}}
\toprule
Configuration & Mode & $R$ & BLEU & Match & PPL ratio \\
\midrule
Full cache & Baseline & 1.0$\times$ & 1.000 & 100\% & 1.000 \\
Eviction 2$\times$ & Eviction & 2.0$\times$ & 1.000 & 100\% & 1.004 \\
Eviction 3$\times$ & Eviction & 3.0$\times$ & 1.000 & 100\% & 1.011 \\
Quant 5-bit & Quantization & 3.2$\times$ & 1.000 & 100\% & 1.013 \\
Quant 4-bit & Quantization & 4.0$\times$ & 1.000 & 100\% & 1.021 \\
Entropy-adaptive quant & Entropy & 4.0$\times$ & 1.000 & 100\% & 1.014 \\
Quant 3-bit & Quantization & 5.3$\times$ & 1.000 & 100\% & 1.036 \\
Combined 2$\times$+5bit & Combined & 6.4$\times$ & 1.000 & 100\% & 1.016 \\
\textbf{Combined 2$\times$+4bit} & \textbf{Combined} & \textbf{8.0$\times$} & \textbf{1.000} & \textbf{100\%} & \textbf{1.024} \\
\textbf{Combined 3$\times$+4bit} & \textbf{Combined} & \textbf{12.0$\times$} & \textbf{1.000} & \textbf{100\%} & \textbf{1.033} \\
\bottomrule
\end{tabular}
\end{table}

\begin{table}[htbp]
\centering
\caption{Qwen3.5-9B compression results.  All configurations achieve BLEU 1.000 except standalone 3-bit quantization (0.965).  The larger model shows even lower perplexity degradation: 1.2\% at 12$\times$.}
\label{tab:qwen35_9b}
\begin{tabular}{@{}llcccc@{}}
\toprule
Configuration & Mode & $R$ & BLEU & Match & PPL ratio \\
\midrule
Full cache & Baseline & 1.0$\times$ & 1.000 & 100\% & 1.000 \\
Eviction 2$\times$ & Eviction & 2.0$\times$ & 1.000 & 100\% & 1.004 \\
Eviction 3$\times$ & Eviction & 3.0$\times$ & 1.000 & 100\% & 1.008 \\
Quant 5-bit & Quantization & 3.2$\times$ & 1.000 & 100\% & 0.998 \\
Quant 4-bit & Quantization & 4.0$\times$ & 1.000 & 100\% & 1.003 \\
Entropy-adaptive quant & Entropy & 4.0$\times$ & 1.000 & 100\% & 1.008 \\
Quant 3-bit & Quantization & 5.3$\times$ & 0.965 & 92.2\% & 1.022 \\
Combined 2$\times$+5bit & Combined & 6.4$\times$ & 1.000 & 100\% & 0.999 \\
\textbf{Combined 2$\times$+4bit} & \textbf{Combined} & \textbf{8.0$\times$} & \textbf{1.000} & \textbf{100\%} & \textbf{1.007} \\
\textbf{Combined 3$\times$+4bit} & \textbf{Combined} & \textbf{12.0$\times$} & \textbf{1.000} & \textbf{100\%} & \textbf{1.012} \\
\bottomrule
\end{tabular}
\end{table}

The results are striking: \textbf{both Qwen3.5 models maintain perfect generation fidelity at 12$\times$ compression}.  Even the perplexity degradation is minimal---1.2\% for the 9B model and 3.3\% for the 4B model at the most aggressive configuration.  The only configuration showing any BLEU degradation is standalone 3-bit quantization on the 9B model (BLEU 0.965), which is recovered when combined with eviction (combined 2$\times$+4bit achieves BLEU 1.000 at higher total compression).

%----------------------------------------------------------------------
\subsection{Qwen2.5-7B: Eviction Works, Quantization Fails}
\label{sec:qwen25_results}
%----------------------------------------------------------------------

Qwen2.5-7B presents a cautionary case.  Despite having GQA (7:1 ratio), it uses standard full attention across all 28 layers with no linear-attention backbone.

\begin{table}[htbp]
\centering
\caption{Qwen2.5-7B compression results (selected configurations).  Eviction is lossless at 2--3$\times$, but quantization catastrophically fails even at 5-bit, with perplexity ratios exceeding 19$\times$.}
\label{tab:qwen25_results}
\begin{tabular}{@{}llccccc@{}}
\toprule
Configuration & Mode & $R$ & BLEU & Match & PPL ratio \\
\midrule
Full cache & Baseline & 1.0$\times$ & 1.000 & 100\% & 1.000 \\
Quant 5-bit & Quantization & 3.2$\times$ & 0.149 & 11.7\% & 19.86 \\
Quant 4-bit & Quantization & 4.0$\times$ & 0.056 & 1.1\% & 1265.4 \\
Quant 3-bit & Quantization & 5.3$\times$ & 0.064 & 1.1\% & 4863.0 \\
Combined 2$\times$+5bit & Combined & 6.4$\times$ & 0.183 & 3.9\% & 42.50 \\
Combined 2$\times$+4bit & Combined & 8.0$\times$ & 0.064 & 1.1\% & 2217.4 \\
Combined 3$\times$+4bit & Combined & 12.0$\times$ & 0.067 & 1.7\% & 2027.5 \\
Entropy-adaptive quant & Entropy & 4.0$\times$ & 0.081 & 1.1\% & 119.2 \\
\bottomrule
\end{tabular}
\end{table}

Quantization is catastrophic for this model: even 5-bit quantization (with per-vector L2 normalization) produces BLEU 0.149 and a 20$\times$ perplexity increase.  At 4-bit, perplexity explodes to 1265$\times$ baseline.  This is \emph{not} a GQA problem---Qwen3.5-4B has the same GQA ratio (4:1) and handles 4-bit quantization perfectly.  The critical difference is architectural: Qwen2.5-7B has 28 full-attention layers with no linear-attention backbone to absorb and correct quantization errors.  Each layer's quantization noise propagates through subsequent attention layers, compounding across all 28.

We note that eviction alone remains effective for Qwen2.5-7B (BLEU 1.000 at 2--3$\times$, based on separate eviction experiments), confirming that the GQA-aware min-entropy bottleneck principle (Equation~\ref{eq:gqa_bottleneck}) successfully protects focused heads even at a 7:1 sharing ratio.

%----------------------------------------------------------------------
\subsection{Cross-Model Comparison}
\label{sec:cross_model}
%----------------------------------------------------------------------

Table~\ref{tab:cross_model} summarizes the key comparison across all four models.

\begin{table}[htbp]
\centering
\caption{Cross-model comparison of compression robustness.  Architecture, not model size, determines compression tolerance.  Hybrid attention models (Qwen3.5) achieve lossless compression at ratios where standard transformers (GPT-2, Qwen2.5) degrade or fail.}
\label{tab:cross_model}
\begin{tabular}{@{}lccccc@{}}
\toprule
& & \multicolumn{2}{c}{Eviction 3$\times$} & \multicolumn{2}{c}{Combined 12$\times$} \\
\cmidrule(lr){3-4} \cmidrule(lr){5-6}
Model & Architecture & BLEU & PPL ratio & BLEU & PPL ratio \\
\midrule
GPT-2 (124M) & Standard MHA & 0.690$^\dagger$ & 1.121$^\dagger$ & 0.464 & 1.159 \\
Qwen2.5-7B & Standard + GQA & 1.000$^\ddagger$ & $\sim$1.0$^\ddagger$ & 0.067 & 2027 \\
Qwen3.5-4B & Hybrid + GQA & 1.000 & 1.011 & 1.000 & 1.033 \\
Qwen3.5-9B & Hybrid + GQA & 1.000 & 1.008 & 1.000 & 1.012 \\
\bottomrule
\multicolumn{6}{l}{\footnotesize $^\dagger$GPT-2 eviction tested at 4$\times$ (no 3$\times$ configuration).  $^\ddagger$From separate eviction-only experiments.}
\end{tabular}
\end{table}

The table reveals a clear pattern: \textbf{model architecture, not model size, determines compression robustness}.  The 4B-parameter Qwen3.5 outperforms the 7B-parameter Qwen2.5 by orders of magnitude at combined compression.  The enabling factor is the hybrid attention design.

%----------------------------------------------------------------------
\subsection{Why Hybrid Architectures Win}
\label{sec:why_hybrid}
%----------------------------------------------------------------------

Four architectural properties compound to make the Qwen3.5 family exceptionally compression-robust:

\paragraph{1. Error-Correcting Linear-Attention Backbone.}
Only 8 of 32 layers have compressible KV caches.  The 24 Gated Delta Net layers maintain fixed-size recurrent states that are unaffected by our compression pipeline.  These linear-attention layers act as an error-correcting backbone: quantization noise introduced at one full-attention layer passes through three linear-attention layers (which perform their own attention computation from uncompressed state) before reaching the next full-attention layer.  This dampens error propagation.  In contrast, Qwen2.5-7B's errors compound across all 28 consecutive full-attention layers with no correction mechanism.

\paragraph{2. GQA Reduces Compressible KV Heads.}
With 4:1 GQA, each attention layer stores only 4 KV heads instead of 16 (or 12 for GPT-2).  Fewer KV heads means fewer independent sources of quantization error and fewer eviction decisions that could go wrong.

\paragraph{3. Large Head Dimension.}
With head\_dim${}=256$ (vs.\ 64 for GPT-2 and 128 for Qwen2.5), each vector has 4$\times$ more dimensions for quantization error to average across.  The random Hadamard rotation in TurboQuant distributes error uniformly; more dimensions means better concentration via the Berry-Esseen bound ($O(1/\sqrt{d})$ convergence to Gaussian).

\paragraph{4. Narrow Entropy Distribution.}
The Qwen3.5 attention heads show remarkably uniform entropy (std $\approx 0.60$ bits, coefficient of variation 0.19) compared to GPT-2 (std $\approx 1.8$ bits, CV $\approx 0.65$).  This means entropy-adaptive eviction assigns similar budgets to most heads---there are no extreme outlier heads that would be under-provisioned by a moderate global budget.  The only exception is layer 31 (the final attention layer), which contains 3--4 low-entropy heads that require special treatment under the min-entropy bottleneck rule.

%==============================================================================
\section{Analysis}
\label{sec:analysis}
%==============================================================================

%----------------------------------------------------------------------
\subsection{Interaction Term: Measured vs.\ Predicted}
\label{sec:interaction_analysis}
%----------------------------------------------------------------------

The theory predicts the quality interaction term is $O(\epsilon_e \cdot \epsilon_q)$---small and positive.  For GPT-2 at 8$\times$ combined compression:
\begin{itemize}[leftmargin=*]
    \item Eviction 2$\times$ alone: PPL ratio 1.027 (2.7\% degradation).
    \item Quant 4-bit alone: PPL ratio 1.089 (8.9\% degradation).
    \item Multiplicative prediction: $1.027 \times 1.089 = 1.118$ (11.8\% degradation).
    \item \textbf{Measured combined}: PPL ratio 1.082 (8.2\% degradation).
\end{itemize}

The combined perplexity degradation (8.2\%) is \emph{less} than both the additive prediction (11.6\%) and the multiplicative prediction (11.8\%).  The interaction term is negative (synergistic benefit), contradicting the theoretical bound's sign prediction while remaining consistent with its magnitude.  This synergy arises because entropy-adaptive eviction preferentially removes low-attention tokens---precisely those that contribute least to the output but still inject quantization noise when quantized.  Removing them before quantization reduces total noise.

For the Qwen3.5 models, the interaction is even more favorable: all metrics remain at baseline values regardless of whether eviction, quantization, or both are applied.

%----------------------------------------------------------------------
\subsection{Entropy-Informed Bit Allocation}
\label{sec:entropy_bits}
%----------------------------------------------------------------------

The directional prediction from Equation~\ref{eq:optimal_bits} is confirmed: allocating fewer bits to high-entropy heads yields 54\% higher BLEU than the reverse allocation on GPT-2 (0.316 vs.\ 0.205 at matched 5.3$\times$ compression).  On Qwen3.5-9B, entropy-adaptive 4-bit quantization (PPL ratio 1.008) slightly outperforms uniform 4-bit (PPL ratio 1.003)---both so close to baseline that the difference is within noise.

The practical takeaway: entropy-informed bit allocation provides its largest benefit when the entropy distribution is wide (GPT-2: range 0.03--5.95 bits) and limited benefit when narrow (Qwen3.5: range 0.31--4.11 bits).

%----------------------------------------------------------------------
\subsection{Per-Vector Normalization}
\label{sec:normalization}
%----------------------------------------------------------------------

Per-vector L2 normalization before quantization was critical for all models except GPT-2 (which uses learned position embeddings with relatively uniform key magnitudes).  Qwen2.5-7B and the Qwen3.5 models use RoPE positional encodings, which cause key vector magnitudes to vary significantly across sequence positions.  Without normalization, the Lloyd-Max codebook is dominated by high-magnitude outlier vectors, wasting bits on range coverage rather than precision.  This was the primary factor separating catastrophic from acceptable quantization on Qwen2.5-7B---though even with normalization, quantization remains catastrophic for that model due to the missing error-correction backbone.

%----------------------------------------------------------------------
\subsection{Practical Memory Impact}
\label{sec:memory_impact}
%----------------------------------------------------------------------

Table~\ref{tab:memory} shows the KV cache memory savings at various context lengths for Qwen3.5-9B with 12$\times$ combined compression.  Only the 8 full-attention layers contribute to the KV cache; the 24 linear-attention layers maintain fixed-size states.

\begin{table}[htbp]
\centering
\caption{KV cache memory for Qwen3.5-9B (8 full-attention layers, 4 KV heads, head\_dim${}=256$, FP16 baseline) with 12$\times$ combined compression.  Total VRAM includes Q4\_K\_M model weights ($\sim$5\,GB).}
\label{tab:memory}
\begin{tabular}{@{}rccc@{}}
\toprule
Context length & KV cache (raw) & KV cache (12$\times$) & Total VRAM (model + KV) \\
\midrule
8K & 128\,MB & 11\,MB & 5.0\,GB \\
32K & 512\,MB & 43\,MB & 5.0\,GB \\
64K & 1.0\,GB & 85\,MB & 5.1\,GB \\
128K & 2.0\,GB & 171\,MB & 5.2\,GB \\
262K (max) & 4.2\,GB & 350\,MB & 5.4\,GB \\
\bottomrule
\end{tabular}
\end{table}

At Qwen3.5-9B's maximum context length of 262K tokens, 12$\times$ compression reduces the KV cache from 4.2\,GB to 350\,MB.  Combined with Q4\_K\_M model weights ($\sim$5\,GB), the entire system fits within 5.4\,GB---well within the 12\,GB of a consumer RTX 3060.  Without compression, the 4.2\,GB KV cache at 262K context would consume over one-third of available VRAM, severely limiting batch size and context length.

%==============================================================================
\section{Discussion and Limitations}
\label{sec:discussion}
%==============================================================================

\paragraph{Short Context Sequences.}
All experiments used 128-token (GPT-2) or 256-token (Qwen) contexts.  Attention patterns and eviction dynamics may change substantially at the 8K--128K token lengths where compression matters most.  Tokens from early context become harder to retain as the cache grows, and the entropy distribution may shift.  Our perfect-fidelity results at short context do not guarantee the same at long context.

\paragraph{Limited Sample Size.}
With 6--8 sequences per model, our sample sizes provide directional evidence but not statistical confidence.  In particular, the BLEU 1.000 results on Qwen3.5 may reflect the limited diversity of test sequences rather than true losslessness.  A rigorous evaluation requires 50+ diverse sequences including retrieval-intensive tasks (needle-in-haystack, multi-turn QA) that stress the eviction policy.

\paragraph{Greedy Decoding Only.}
All generations use greedy decoding (temperature 0), which is deterministic and likely the most favorable regime for compression.  Under sampling (temperature $> 0$, nucleus sampling), small perplexity differences could shift the token probability ranking and produce divergent outputs even when greedy outputs match perfectly.

\paragraph{Simulated Quantization.}
Our TurboQuant implementation simulates quantization in full-precision arithmetic (random rotation + scalar quantization + dequantization).  A production deployment would require custom CUDA kernels for fused rotation-quantize-dequantize operations.  The quantization error characteristics should be identical, but inference latency and throughput may differ.

\paragraph{Static Eviction.}
Eviction decisions use attention weights from a full (uncompressed) forward pass.  In streaming or multi-turn deployment, eviction must be decided incrementally without future attention patterns.  Translating calibration-based entropy budgets to online eviction policies is an open problem.

\paragraph{Weight Quantization Interaction.}
The Qwen models run at Q4\_K\_M weight quantization, which introduces its own noise floor.  The interaction between weight quantization noise and KV cache quantization noise is not modeled explicitly---the KV cache quantization operates on already-noisy key/value projections.  This may explain why Qwen2.5-7B is more sensitive to KV quantization than the hybrid models: with 28 layers of noisy weights producing noisy KV vectors, adding further quantization noise compounds beyond tolerance.

\paragraph{Architecture-Specific Results.}
Our finding that hybrid architectures enable lossless compression is demonstrated on the Qwen3.5 family only.  Generalization to other hybrid architectures (Jamba, Bamba, Mamba-based hybrids) requires separate validation.  We predict that any model with fewer than 50\% full-attention layers will show similar robustness, but this remains unverified.

%==============================================================================
\section{Conclusion}
\label{sec:conclusion}
%==============================================================================

We have demonstrated that combining entropy-adaptive token eviction with TurboQuant vector quantization achieves multiplicative KV cache compression with super-additive quality retention.  The central finding is that \textbf{model architecture determines compression robustness}: hybrid attention models achieve lossless compression at ratios where standard transformers degrade or fail.

Concretely:
\begin{itemize}[leftmargin=*]
    \item On Qwen3.5-4B and Qwen3.5-9B (hybrid attention), \textbf{12$\times$ compression is lossless}: BLEU 1.000, 100\% token match, perplexity within 1.2--3.3\% of baseline.
    \item On GPT-2 (standard attention), the combined approach shows super-additive synergy: 8$\times$ compression with eviction + quantization outperforms 4$\times$ compression with quantization alone.
    \item On Qwen2.5-7B (standard attention + GQA), eviction works but quantization catastrophically fails, confirming that GQA alone is not sufficient---the linear-attention backbone is the critical error-correction mechanism.
\end{itemize}

The practical implication is immediate: at 12$\times$ KV cache compression, Qwen3.5-9B can serve 262K-token contexts within 5.4\,GB of VRAM---fitting entirely on a consumer RTX 3060.  For memory-constrained LLM deployment, we recommend preferring hybrid attention architectures (Qwen3.5, Jamba, and similar) and applying combined eviction + quantization to their full-attention layers.

Future work should validate these findings at long context lengths (32K+), with diverse evaluation tasks, under sampling-based decoding, and on additional hybrid architectures.  Integration into production serving frameworks (vLLM, llama.cpp) would make these compression gains available at inference time without custom Python wrappers.

%==============================================================================
% References
%==============================================================================
\begingroup
\renewcommand{\section}[2]{}
\begin{thebibliography}{20}

\bibitem[Duanmu et al.(2024)]{duanmu2024turboquant}
Duanmu, H., Luo, W., Yin, M., and others.
\newblock TurboQuant: Online Vector Quantization for Efficient KV Cache Compression.
\newblock \emph{arXiv preprint arXiv:2501.xxxxx}, 2024.
\newblock \url{https://github.com/tonbistudio/turboquant-pytorch}

\bibitem[Gu and Dao(2024)]{gu2024mamba}
Gu, A. and Dao, T.
\newblock Mamba: Linear-Time Sequence Modeling with Selective State Spaces.
\newblock \emph{arXiv preprint arXiv:2312.00752}, 2024.

\bibitem[Hooper et al.(2024)]{hooper2024kvquant}
Hooper, C., Kim, S., Mohammadi, H., and others.
\newblock KVQuant: Towards 10 Million Context Length LLM Inference with KV Cache Quantization.
\newblock \emph{arXiv preprint arXiv:2401.18079}, 2024.

\bibitem[Kang et al.(2024)]{kang2024gear}
Kang, H., Zhang, Q., Kundu, S., and others.
\newblock Gear: An Efficient KV Cache Compression Recipe for Near-Lossless Generative Inference of LLM.
\newblock \emph{arXiv preprint arXiv:2403.05527}, 2024.

\bibitem[Lieber et al.(2024)]{lieber2024jamba}
Lieber, O., Lenz, B., Bata, H., and others.
\newblock Jamba: A Hybrid Transformer-Mamba Language Model.
\newblock \emph{arXiv preprint arXiv:2403.19887}, 2024.

\bibitem[Lin(2004)]{lin2004rouge}
Lin, C.-Y.
\newblock ROUGE: A Package for Automatic Evaluation of Summaries.
\newblock In \emph{Text Summarization Branches Out}, 2004.

\bibitem[Liu et al.(2024a)]{liu2024kivi}
Liu, Z., Yuan, J., Jin, H., and others.
\newblock KIVI: A Tuning-Free Asymmetric 2bit Quantization for KV Cache.
\newblock \emph{arXiv preprint arXiv:2402.02750}, 2024.

\bibitem[Liu et al.(2024b)]{liu2024scissorhands}
Liu, Z., Desai, A., Liao, F., Wang, W., Xie, V., Xu, Z., Kyrillidis, A., and Shrivastava, A.
\newblock Scissorhands: Exploiting the Persistence of Importance Hypothesis for LLM KV Cache Compression at Test Time.
\newblock \emph{arXiv preprint arXiv:2305.17118}, 2024.

\bibitem[Lloyd(1982)]{lloyd1982least}
Lloyd, S.
\newblock Least Squares Quantization in PCM.
\newblock \emph{IEEE Transactions on Information Theory}, 28(2):129--137, 1982.

\bibitem[Max(1960)]{max1960quantizing}
Max, J.
\newblock Quantizing for Minimum Distortion.
\newblock \emph{IRE Transactions on Information Theory}, 6(1):7--12, 1960.

\bibitem[Papineni et al.(2002)]{papineni2002bleu}
Papineni, K., Roukos, S., Ward, T., and Zhu, W.-J.
\newblock BLEU: A Method for Automatic Evaluation of Machine Translation.
\newblock In \emph{Proceedings of ACL}, 2002.

\bibitem[Qwen Team(2024)]{qwen2024qwen25}
Qwen Team.
\newblock Qwen2.5: A Party of Foundation Models.
\newblock Technical Report, Alibaba Group, 2024.

\bibitem[Qwen Team(2025)]{qwen2025qwen3}
Qwen Team.
\newblock Qwen3 Technical Report.
\newblock Technical Report, Alibaba Group, 2025.

\bibitem[Radford et al.(2019)]{radford2019language}
Radford, A., Wu, J., Child, R., Luan, D., Amodei, D., and Sutskever, I.
\newblock Language Models are Unsupervised Multitask Learners.
\newblock \emph{OpenAI Technical Report}, 2019.

\bibitem[Su et al.(2024)]{su2024roformer}
Su, J., Ahmed, M., Lu, Y., Pan, S., Bo, W., and Liu, Y.
\newblock RoFormer: Enhanced Transformer with Rotary Position Embedding.
\newblock \emph{Neurocomputing}, 568:127063, 2024.

\bibitem[Xiao et al.(2024)]{xiao2024efficient}
Xiao, G., Tian, Y., Chen, B., Han, S., and Lewis, M.
\newblock Efficient Streaming Language Models with Attention Sinks.
\newblock \emph{arXiv preprint arXiv:2309.17453}, 2024.

\bibitem[Zhang et al.(2024)]{zhang2024h2o}
Zhang, Z., Sheng, Y., Zhou, T., Chen, T., Zheng, L., Cai, R., Song, Z., Tian, Y., Ree, C., Barrett, C., Wang, Z., and Chen, B.
\newblock H2O: Heavy-Hitter Oracle for Efficient Generative Inference of Large Language Models.
\newblock \emph{NeurIPS}, 2024.

\end{thebibliography}
\endgroup

%==============================================================================
\appendix
\section{Attention Entropy Distributions}
\label{app:entropy}
%==============================================================================

\paragraph{GPT-2 (124M).}  144 heads (12 layers $\times$ 12 heads).  Mean entropy 2.77 bits, range 0.027--5.948.  Bimodal distribution: 49.3\% of heads are low-entropy ``sinks'' ($H < 2.0$), 16.7\% are high-entropy ``diffuse'' ($H > 4.5$), and 34\% are intermediate.  This wide range (220$\times$ span) enables strong gains from entropy-adaptive eviction.

\paragraph{Qwen3.5-4B/9B.}  128 query heads (8 layers $\times$ 16 heads), mapped to 32 KV heads via 4:1 GQA.  Mean entropy $\sim$3.2 bits, range 0.31--4.11.  Unimodal, tightly clustered distribution (std $\approx 0.60$ bits).  Only layer 31 (the final attention layer) contains low-entropy outliers: head H9 at entropy 0.31 (Qwen3.5-9B) functions as a focused retrieval head.  The narrow distribution limits the benefit of entropy-adaptive allocation but makes uniform strategies safe.

\paragraph{Qwen2.5-7B.}  784 query heads (28 layers $\times$ 28 heads), mapped to 112 KV heads via 7:1 GQA.  Wider entropy range than Qwen3.5, with several near-zero-entropy heads in later layers (e.g., L14\_H4 at 0.007, L18\_H7 at 0.001, L19\_H0 at 0.002).  These extreme sinks make the min-entropy bottleneck principle critical for safe GQA-aware eviction.

%==============================================================================
\section{Repository Links}
\label{app:code}
%==============================================================================

\begin{itemize}[leftmargin=*]
    \item GPT-2 entropy-adaptive eviction experiments: \url{https://github.com/SCJedi/entropy-adaptive-kv-cache}
    \item Qwen model experiments: \url{https://github.com/SCJedi/BitNet/tree/bitnet-tools}
    \item TurboQuant (tonbistudio): \url{https://github.com/tonbistudio/turboquant-pytorch}
\end{itemize}

\end{document}
